% บทคัดย่อภาษาไทย
% ย่อหน้าที่ 2 สรุปผลให้ตรงกับบทที่ 4 และ 5 (ฉบับสมบูรณ์ ตุลาคม 2569)
% ส่วนหัว (ชื่อผู้จัดทำ/ปี/ที่ปรึกษา) และคำสำคัญท้ายบทคัดย่อ สร้างอัตโนมัติจาก main.tex ไม่ต้องใส่ในไฟล์นี้

โครงงานนี้มีวัตถุประสงค์เพื่อพัฒนาเว็บแอปพลิเคชัน ``ไดโน'' (Dino) ให้เป็นแพลตฟอร์มศูนย์กลางข้อมูลการท่องเที่ยวในจังหวัดขอนแก่น เพื่อแก้ไขปัญหาการเข้าถึงข้อมูลการท่องเที่ยวที่กระจัดกระจายและไม่ครบถ้วน โครงงานนี้ได้บูรณาการเทคโนโลยีปัญญาประดิษฐ์ (AI) ผ่านโมเดลภาษาขนาดใหญ่ (LLMs) ร่วมกับสถาปัตยกรรม Retrieval-Augmented Generation (RAG) ในการพัฒนาระบบหลัก 2 ส่วน ได้แก่ (1) ระบบสร้างแผนการเดินทางแบบปรับตัว (Adaptive Trip Planning) ที่ประเมินจากข้อจำกัดด้านงบประมาณ ระยะเวลา และความสนใจเฉพาะบุคคล พร้อมกลไกตรวจสอบคุณภาพแผนด้วยเทคนิค LLM-as-a-Judge และ (2) ระบบแชตบอต (Chatbot) สำหรับโต้ตอบและให้คำแนะนำด้านการท่องเที่ยวแบบเรียลไทม์ โดยใช้ฐานข้อมูลสถานที่และร้านอาหารจาก Google Places API ภายในจังหวัดขอนแก่น

ผลการดำเนินงานพบว่า ส่วนหลังบ้าน ส่วนหน้าบ้าน และส่วนบริการปัญญาประดิษฐ์เชื่อมต่อและทำงานร่วมกันได้แบบครบวงจร (End-to-End) และนำขึ้นใช้งานจริงที่ dinokhonkaen.app ทั้งแชตบอตและระบบวางแผนการเดินทางด้วยโมเดลภาษาขนาดใหญ่ รวมถึงระบบบัญชีผู้ใช้ ระบบสแกน QR สะสม/แลกพอยท์ และระบบจัดการหลังบ้าน การทดลองเปรียบเทียบ Prompt ของระบบวางแผนการเดินทางด้วยโมเดลจริง 96 รัน พบว่าแผนผ่านตัวชี้วัดด้านความถูกต้องของสถานที่ สถานที่ที่ผู้ใช้กำหนด และความสมเหตุสมผลของแผน และผลการประเมินความพึงพอใจด้วยแบบสอบถามจากผู้ใช้ \surveyN{} คนได้ค่าเฉลี่ยรวม \surveyOverallMean{} จาก 5 คะแนน อยู่ในระดับ\surveyOverallLevel{} โดยด้านแชตบอตและด้านการวางแผนเที่ยวด้วย AI ได้คะแนนต่ำกว่าด้านอื่น โดยเฉพาะความเร็วในการตอบสนอง ผลดังกล่าวแสดงว่าแอปพลิเคชันสามารถรวมข้อมูลการท่องเที่ยวของจังหวัดขอนแก่นไว้ในแหล่งเดียว และสร้างแผนการเดินทางที่ตรงกับความต้องการเฉพาะบุคคลได้จริง ซึ่งเป็นพื้นฐานในการสนับสนุนผู้ประกอบการท้องถิ่นและเศรษฐกิจการท่องเที่ยวของจังหวัดขอนแก่นต่อไป
